\begin{center}
\begin{tikzpicture}[>=stealth]
 \draw[->] (-0.6,0) -- (9.2,0) node[below] {$x$};
 \draw[->] (0,-0.3) -- (0,3.2) node[left] {$y$};
 \draw[thick, domain=0:8.145, samples=60] plot (\x, {0.85+0.9004*\x-0.12336*\x*\x});
 \fill (0,0.85) circle (0.1);
 \draw[->, very thick, red] (0,0.85) -- ++(42:1.8) node[above left] {$\vec v_0$};
 \draw[dashed] (0,0.85) -- (1.4,0.85);
 \draw (1.0,0.85) arc (0:42:1.0);
 \node at (1.3,1.15) {$\alpha$};
 \draw[|<->|] (-0.4,0) -- (-0.4,0.85) node[midway, left] {$h$};
 \draw[|<->|] (0,-0.5) -- (8.145,-0.5) node[midway, below] {$w$};
\end{tikzpicture}
\end{center}

Die Teilgeschwindigkeiten beim Abstoss sind
\begin{align}
 v_x &= \vxkF = \vk\cdot\cos\left(\alO\right) = \vxk \\
 v_y &= \vykF = \vk\cdot\sin\left(\alO\right) = \vyk
\end{align}
Die Höhe der Kugel über dem Boden ist (nach oben positiv)
\begin{align}
 y(t) &= h + v_y\,t - \frac{1}{2}g\,t^2
\end{align}
Beim Aufprall ist $y = 0$. Die positive Lösung dieser quadratischen Gleichung ist die Flugzeit
\begin{align}
 t &= \tkF \\
   &= \frac{\vyk+\sqrt{\left(\vyk\right)^2+2\cdot\gE\cdot\hk}}{\gE} \\
   &= \tk
\end{align}
Horizontal bewegt sich die Kugel gleichförmig:
\begin{align}
 w &= \xkF = \vxk\cdot\tk = \xk \approx \result{\xkP{0}{2}}
\end{align}
